Para executar todo o ambiente de desenvolvimento localmente, adotou-se o \textit{Docker}, que fornece um ambiente padronizado com \textit{hot reload}, permitindo que as alterações no código sejam refletidas imediatamente, sem necessidade de reiniciar a aplicação.

Como a execução do \textit{frontend} depende do \textit{backend}, optou-se por publicar a versão compilada do \textit{backend} no \textit{Docker Hub}. Dessa forma, quem desenvolve o \textit{frontend} não precisa compilar nem configurar o \textit{backend} em sua máquina. Para integrar os serviços, foi implementado um arquivo \textit{Docker Compose} que cria uma rede privada, na qual o \textit{backend} e o banco de dados local se comunicam de forma isolada do ambiente externo.

A atualização da imagem é automatizada, a cada \textit{commit} na \textit{branch} \texttt{main} do repositório do \textit{backend}, o \textit{Docker Hub} é acionado para compilar uma nova imagem e publicá-la no repositório, garantindo que a versão mais recente esteja sempre disponível para o \textit{frontend}. A imagem mais atualizada do \textit{backend} pode ser acessada em: \url{https://hub.docker.com/r/joaoseisei1/tcc-seed-service}.